% !TeX spellcheck = en_US
% !TeX root = tikz-ext-manual.tex
% Copyright 2026 by Qrrbrbirlbel
%
% This file may be distributed and/or modified
%
% 1. under the LaTeX Project Public License 1.3c and/or
% 2. under the GNU Free Documentation License.
%
\section{Shape: Any}
\begin{purepgflibrary}{ext.shapes.any}
  A shape that can have \emph{any} background path.
\end{purepgflibrary}
\begin{ext_shape}{any}
  This shape is a shape that can have a custom background path and a custom border path.
  For the border path, the \referenceLibraryandIndexO{intersections} library is needed and is loaded automatically.
  If many points on the border are requested this might be quite the bottleneck for your diagrams.
  
  \begin{key}{/\pgfext/any path=\meta{background path} (initially \normalfont{empty})}
    This sets the path for the background of the |any| shape.
    
    Remember that if this stays empty or is smaller than the text of the node,
    the bounding box will not necessarily include the whole text of the node.
    
    The node's internal coordinate system is setup in such a way
    that the origin is at the center of the node's text box.
  \end{key}
  \begin{key}{/\pgfext/any border=\meta{border path} (initially \normalfont{empty})}
    This sets the path for the border of the |any| shape.

    For example, the following example would replicate the background path of the |rectangle| shape.
    
    {\small
      (We'll explicitly use the keys in the |/pgf/ext| namespace here
      since the keys in the |/tikz/ext| namespace will wrap their arguments in \tikzname's path parser.)\par}
\begin{codeexample}[preamble=\usepgflibrary{ext.shapes.any}]
\begin{tikzpicture}[ultra thick]
\node[
  shape=ext_any,
  draw,
  /pgf/ext/any path={
    \pgfpathrectanglecorners
      {\pgfqpoint{ \halftextwidthinner}{ \halftexttotalheightinner}}
      {\pgfqpoint{-\halftextwidthinner}{-\halftexttotalheightinner}}
  },
  /pgf/ext/any border={
    \pgfpathrectanglecorners
      {\pgfpoint{ \halftextwidthinner+\outerxsep}{ \halftexttotalheightinner+\outerysep}}
      {\pgfpoint{-\halftextwidthinner-\outerxsep}{-\halftexttotalheightinner-\outerysep}}
  }
] (fb) {Foo Bar};
\node[draw] at (2, .5) {Foo Baz} edge[red] (fb);
\end{tikzpicture}
\end{codeexample}

    Internally, the \referenceLibraryandIndexO{intersections} library will be used to
    find a point where the border path will intersect with a straight ray coming from the center.
    For efficiency reasons the following values should be set so that this ray will
    be as short as possible.
    \begin{key}{/\pgfext/any min ray radius=\meta{dimension} (initially 0pt)}
      This value defines that the start point of that ray will be \meta{dimension} away from the |center| anchor.
    \end{key}
    \begin{key}{/\pgfext/any max ray radius=\meta{dimension} (initially 100pt)}
      This value defines that the end point of that ray will be \meta{dimension} away from the |center| anchor.
    \end{key}
  \end{key}

  If no |border path| is set and the \referenceLibraryandIndexExt{offsetpath} library \cite{nfold}
  is loaded, it can be used to automatically offset the |any path| according to the following key to
  find the border path.
  \begin{key}{/\pgfext/any use offset=\mchoice{false, right, left, counter clockwise, clockwise, ccw, cw} (default false)}
    This key disables (|false|) or enables (all other options) the use of the |offsetpath| library to
    automatically offset the |any path|.
    The choices |right|, |counter clockwise| and |ccw| are aliases of each other,
    as are |left|, |clockwise| and |cw|.
    
    The distance for offset is the maximum of both |outer sep|s.
    
    This makes it very convenient to create a border for complex \enquote{shapes} but
    remember that the use of |offsetpath| neither is efficient nor can it guarantee perfect results.
\begin{codeexample}[preamble=\usepgflibrary{ext.shapes.any, offsetpath}]
\begin{tikzpicture}[ultra thick]
\node[
  shape=ext_any,
  draw,
  /pgf/ext/any path={
    \pgfpathrectanglecorners
      {\pgfqpoint{ \halftextwidthinner}{ \halftexttotalheightinner}}
      {\pgfqpoint{-\halftextwidthinner}{-\halftexttotalheightinner}}
  },
  ext/any use offset=left
] (fb) {Foo Bar};
\node[draw] at (2, .5) {Foo Baz} edge[red] (fb);
\end{tikzpicture}
\end{codeexample}
  \end{key}

Of course, no shape would be complete without named anchors,
for this task, the following keys are provided.

\begin{key}{/\pgfext/any add anchor=\marg{name}\marg{coordinate}}
  This will add the anchor named \meta{name} at the \meta{coordinate} (remember, this is from the text's center).
\end{key}
\begin{key}{/\pgfext/any add direction anchor=\marg{name}\marg{direction}}
  Similar to above, but it will use the shape's border to find the a point in the given \meta{direction} (a number in degree).
\end{key}

\clearpage

As you've noticed, various macros can be used inside the setup of the node.
Any of these may appear inside of |any min ray radius|, |any max ray radius|
or an anchor specification.
\begin{commandlist}{%
  \minimumwidth, \minimumheight, \innerxsep, \innerysep, \outerxsep, \outerysep%
}
These commands return the evaluated values of the corresponding |/pgf| value-keys.
\begin{commandlist}{%
  \minimumsize, \innersep, \outersep
}
Some shapes only take the maximum of the previous values in consideration,
these commands serve as a short-cut to these.
\end{commandlist}
\end{commandlist}

\begin{commandlist}{%
  \halftextwidth,
  \halftexttotalheight,
  \halftextwidthinner,
  \halftexttotalheightinner
}
These commands return lengths corresponding to the text box. The ones with suffix |inner| include either |inner xsep| or |inner ysep|.

As the previous example shows, the last two correspond to (half the) dimensions of a |rectangle| node.
The first two are just neat wrappers for (half the) size of \referenceCommandandIndexO{\pgfnodeparttextbox}
(which is also available, of course).
\end{commandlist}

%\clearpage

\begin{tikzlibrary}{ext.shapes.any}
  A shape that can have \emph{any} background path.
\end{tikzlibrary}

The \tikzname\ library offers a few wrappers for using the shape with \tikzname.
\begin{key}{/\tikzext/any path=\meta{background path} (initially \normalfont{empty})}
  As |/pgf/ext/any path| but \tikzname\ path specifications are supported.
\end{key}
\begin{key}{/\tikzext/any border=\meta{background path} (initially \normalfont{empty})}
  As |/pgf/ext/any border| but \tikzname\ path specifications are supported.
\end{key}
\begin{key}{/\tikzext/any add anchor=\meta{name}\meta{coordinate}}
  As |/pgf/ext/any add anchor| but \tikzname\ coordinate specifications (without |(| and |)|) are supported.
\end{key}

\clearpage
\begin{codeexample}[preamble=\usetikzlibrary{ext.any, offset},width=16cm]
\begin{tikzpicture}\Huge
\node[name=s, shape=ext_any, shape example, line join=round,
  ext/any path={
       ( \halftextwidthinner,  \halftexttotalheightinner) -- ( 0pt                           ,  \halftexttotalheightinner+\innerysep)
    -- (-\halftextwidthinner,  \halftexttotalheightinner) -- (-\halftextwidthinner-\innerxsep, 0pt                                  )
    -- (-\halftextwidthinner, -\halftexttotalheightinner) -- ( 0pt                           , -\halftexttotalheightinner-\innerysep)
    -- ( \halftextwidthinner, -\halftexttotalheightinner) -- ( \halftextwidthinner+\innerxsep, 0pt                                  )
    -- cycle
  },
  ext/any use offset=right,
  ext/any add directional anchor/.list={
    {east}{0}, {north}{90}, {west}{180}, {south}{270}, {north east}{45}, {north west}{135}, {south west}{225}, {south east}{315}},
  ext/any min ray radius={min(\halftextwidth,\halftexttotalheight)-1pt},
  ext/any max ray radius={max(\halftextwidthinner,\halftexttotalheightinner)+\innersep+5*\outersep+1pt},
]
  {Any Shape\vrule width 1pt height 2cm};
\foreach \anchor/\placement in {
  north west/above left, north/above, north east/above right, west/left, center/above, east/right,
  mid/above, base/below, % mid west/right, mid east/left, base west/left, base east/right,
  south west/below left, south/below, south east/below right, text/left, 10/right, 130/above}
  \draw[shift=(s.\anchor)] plot[mark=x] coordinates{(0,0)}
    node[\placement] {\scriptsize\texttt{(s.\anchor)}};
\end{tikzpicture}
\end{codeexample}
\end{ext_shape}
\endinput